\section{Conclusion}

QKF now connects finite algebraic consequences to two executable proof paths: a checked source interface that can be reused by consumers, and a direct source-to-consumer equality proof from scoped native facts. The first retains labeled observational distinctions and composes their actions across words. The second records exactly how native premises and earlier closed lemmas support the requested result. Initial pure-row semantics supplies a separate exact account of forced values, including values outside the generated domain.

The original coordinatewise characterization and signed-bound contracts remain intact. Later caller profiles and numeric target theorems strengthen particular composition steps without claiming formal verification of arbitrary Java. Lean now also establishes soundness of decoded typed positive ground proofs and conservative chronological lemma import; source semantics, adapters, and execution bridges retain explicit obligations.

The measurements give a concrete engineering direction. Prepared contexts and direct emission remove significant repeated work, while QKF's simpler ordinary-SDK route remains the stronger registered cost baseline. This supports continued development of useful checked summaries and disciplined selection among internal backends. Wider source coverage, richer observation languages, and an advantage attributable to the more elaborate query and forcing mechanisms remain questions for controlled experiments and further proofs.
